% Source: https://artofproblemsolving.com/community/c3h1980862_amc_10_question
% Stable identity: aops-topic-1980862-q01
\emph{Modes.} The values $1,2,\dots,28$ each occur $12$ times, more than any other value, so they are the modes. Their median is $d=\dfrac{14+15}{2}=14.5$.

\emph{Median.} The median of $365$ values is the $183$rd smallest. The values $1$ to $15$ account for $12\cdot15=180$ data, and the $16$s occupy places $181$ to $192$. So $M=16$.

\emph{Mean, without a long computation.} If every date $1,\dots,31$ occurred $12$ times, the $372$ values would be symmetric about $16$ and the mean would be $16$. The actual data remove one $29$, one $30$ and five $31$s: seven values, all larger than $16$, so the mean drops below $16$, i.e.\ $\mu<M$. It does not drop much: the total becomes $372\cdot16-(29+30+5\cdot31)=5952-214=5738$, and
\[
\mu=\dfrac{5738}{365}\approx15.72>14.5=d .
\]
(Even more simply, $14.5\cdot365=5292.5<5738$.) Hence $d<\mu<M$, answer $\boxed{\textbf{(E)}}$.
